
\section{收缩 $(-1)$-\trCurve{}的\trCastelnuovo{}判别准则}

我们在\mbox{练习~22.4.O} 中已经看到：在一张\trRegular\trSurface{}的一个（既约）点处作\trBlowup，所得的\trExceptionalDivisor{}是一条 $(-1)$-\trCurve{}（见\mbox{定义~20.2.8}：它\trIsomorphic{}于 $\mathbb{P}^1$，且其\trNormalBundle{}为 $\OO_C(-1)$）。\trCastelnuovo{}判别准则是它的逆命题：若我们有一张含有 $(-1)$-\trCurve{}的\trQuasiprojective{}\trSurface，那么这张曲面就是在另一张曲面的一个既约正则点处作\trBlowup{}得到的。我们就说这条 $(-1)$-\trCurve{}可以被\textbf{收缩}（blow down）。
\textbf{28.6.1. \trhTheorem（\trCastelnuovo{}判别准则）。}\ ——\ \emph{设 $C\subset X$ 是 $k$ 上一张\trSmooth\trProjective{}\trSurface{}上的一条 $(-1)$-\trCurve。则存在一个\trBirational\trMorphism{} $\pi\colon X\to Y$，使得 $Y$ 是一张\trSmooth\trProjective{}\trSurface，$\pi(C)$ 是一个 $k$-\trValuedPoint，$\pi$ 是 $Y$ 在该点处的\trBlowup，而 $C$ 是这次\trBlowup{}的\trExceptionalDivisor。}
由\mbox{练习~28.4.C}，「收缩 $C$」的办法只有一种：这个收缩是唯一的。
\emph{证明。}\quad 证明分三步。\emph{第 1 步.} 我们构造 $\pi\colon X\to Y$，它把 $C$ 收缩为一点 $q$（且在其余之处是\trIsomorphism）。到这一步完成时，我们将知道 $Y$ 是一个\trProjective\trVariety。\emph{第 2 步.} 然后我们证明 $Y$ 是\trSmooth{}的；这是最艰难的一步，需要用到\mbox{形式函数定理~28.3.2}。\emph{第 3 步.} 最后，我们认出 $\pi$ 就是\trBlowup{} $\mathrm{Bl}_q\,Y\to Y$。
\emph{第 1 步.}\quad 我们的思路是找出一个\trInvertibleSheaf{} $\mathscr{L}$，让它的完全线性系来承担这项任务。我们从 $X$ 上一个满足 $\mathrm{H}^1(X,\mathscr{H})=0$ 的\trVeryAmple\trInvertibleSheaf{} $\mathscr{H}$ 出发。这样的\trInvertibleSheaf{} $\mathscr{H}$ 可以如下找到：任取一个\trVeryAmple\trInvertibleSheaf，再取其一个适当大的倍数，其间援引\mbox{练习~16.2.C}（\trVeryAmple{} $\otimes$ \trVeryAmple{} $=$ \trVeryAmple{}）与\trSerre{}消失（\mbox{定理~18.1.3}~(ii)）。令 $d=\mathscr{H}\cdot C:=\deg_C\mathscr{H}|_C$。由 $\mathscr{H}$ 的\trVeryAmpleness{}（\mbox{练习~20.1.K}）可知 $d>0$。
\textbf{28.6.A.（练习）}\quad 证明：$\mathrm{H}^1(X,\mathscr{H}(iC))=0$（$0\le i\le d$）。提示：对 $i$ 作归纳，利用
\begin{equation*}
\tag{28.6.1.1}\label{eq:castelnuovo-restriction-seq}
0\longrightarrow\mathscr{H}((i-1)C)\longrightarrow\mathscr{H}(iC)\longrightarrow\mathscr{H}(iC)|_C\longrightarrow 0.
\end{equation*}
提示：注意 $\mathscr{H}(iC)|_C$ 是 $C\cong\mathbb{P}^1$ 上的一个\trLineBundle；它是哪一个 $\OO(n)$？
令 $\mathscr{L}:=\mathscr{H}(dC)$。我们将经由映射 $\phi\colon\mathrm{H}^0(X,\mathscr{H})\to\mathrm{H}^0(X,\mathscr{H}(dC))$，用 $\mathscr{H}$ 的\trSection{}来得到 $\mathscr{L}$ 的\trSection。
\textbf{28.6.B.（练习）}\quad 证明：$\mathscr{L}$ 是\trBasePointFree{}的。提示：要证明 $\mathscr{L}$ 在 $C$ 之外没有基点，考虑 $\phi$ 的像。要证明 $\mathscr{L}$ 在 $C$ \emph{上}没有基点，则利用：由 \eqref{eq:castelnuovo-restriction-seq}，当 $i=d$ 时，$\mathrm{H}^0(X,\mathscr{L})\to\mathrm{H}^0(C,\mathscr{L}|_C)$ 是满射，并且 $\mathscr{L}|_C\cong\OO_C$。
于是，完全线性系 $|\mathscr{L}|$ 给出一个从 $X$ 到某个\trProjectiveSpace{} $\mathbb{P}^N$ 的\trMorphism{} $\pi'$。
\textbf{28.6.C.（练习）}\quad 证明：$\pi'$ 恰好把 $C$ 收缩为一点 $q'$。更明确地说，证明：$\pi'$ 把 $C$ 送到一点 $q'$，而 $\pi'|_{X\setminus C}\colon X\setminus C\to\mathbb{P}^N\setminus\{q'\}$ 是一个\trClosedEmbedding。提示：要证明 $\pi'$ 在 $C$ 之外给出一个\trClosedEmbedding，注意 $|\mathscr{H}|$ 分离点与切向量，并论证 $|\mathscr{L}|=|\mathscr{H}(dC)|$ 在 $C$ 之外分离点与切向量。要证明 $C$ 被 $\pi'$ 收缩，则利用 $\deg\mathscr{L}|_C=0$ 这一事实。
令 $Y'$ 为 $\pi'$ 的像，于是我们有一个\trMorphism{} $X\to Y'$（我们也称它为 $\pi'$）。令 $\nu\colon Y\to Y'$ 为 $Y'$ 的\trNormalization，于是由正规化的\trUniversalProperty{}（\S10.7），$\pi'$ 提升为 $\pi\colon X\to Y$。（我们作正规化，是为了能在第 2 步中用上\mbox{练习~28.4.B}。）
\textbf{28.6.D.（容易练习）}\quad 证明：$\nu^{-1}(q')\subset Y$ 是单独一个点，我们把它称作 $q$，并且 $\pi(C)=q$。（提示：$\pi$ 的像是\trClosed{}的，而 $\pi(C)$ 是\trConnected{}的。）
\textbf{28.6.E.（练习）}\quad 设 $X$ 是一张\trSmooth\trProjective{}\trSurface，含有一条\trIsomorphic{}于 $\mathbb{P}^1_k$ 的\trCurve{} $C\subset X$。若 $C\cdot C:=\deg_C\OO(C)|_C<0$，证明：存在一个\trMorphism{} $\pi\colon X\to X'$，它把 $C$ 收缩为一点，而让 $X$ 的其余部分保持不变。（这条练习不会被用到，收在这里是为了让人练习第 1 步中的论证。）
\emph{第 2 步.}\quad 现在我们来证明 $Y$ 是一张\trSmooth\trSurface。由于 $Y$ 是\trNormal{}的，由\mbox{练习~28.4.B}，我们有 $\pi_*\OO_X\cong\OO_Y$（即 $\pi\colon X\to Y$ 是 $\OO$-连通的）。我们将这样来证明 $Y$ 在 $q$ 处\trSmooth：证明作为拓扑环有 $\hat{\OO}_{Y,q}\cong k[[x,y]]$，因为这样一来便有 $\dim_k(\mathfrak{m}/\mathfrak{m}^2)=2$。（你明白为什么这就表明 $Y$ 是\trSmooth{}的，而不仅仅是\trRegular{}的吗？）由\mbox{形式函数定理~28.3.2}，$\hat{\OO}_{Y,q}\cong\lim_{n\to\infty}\mathrm{H}^0(C'_n,\OO_{C'_n})$，其中 $C'_n$ 是 $X$ 的这样一个\trClosedSubscheme：它是 $q$ 的 $n$ 阶形式邻域的\trSchemeTheoreticPreimage。
我们并不确切知道这个 $n$ 阶形式邻域是什么，但这无妨：这个逆向系统与 $\OO(-nC)=\mathscr{I}^n_{C/X}$ 共尾（你看出是为什么吗？），所以我们可以改取后者的\trInverseLimit{}（据\mbox{练习~28.1.A}）。只须证明：对一切 $n\ge 0$ 都有 $\mathrm{H}^0(C_n,\OO_{C_n})\cong k[[x,y]]/(x,y)^n$，其中 $C_n$ 由下式定义
\[
0\longrightarrow\OO(-nC)\longrightarrow\OO_X\longrightarrow\OO_{C_n}\longrightarrow 0.
\]
（当心：我们并不知道 $C_n=C'_n$，甚至连 $C_1=C'_1$ 也不知道。就我们所知，$q$ 之上的\trFiber{}完全可能是非既约的！）
我们按 $n$ 归纳来证。作为 $k$-\trVectorSpace{}这很容易，只需利用
\begin{equation*}
\tag{28.6.1.2}\label{eq:castelnuovo-graded-piece-seq}
0\longrightarrow\mathscr{I}^n/\mathscr{I}^{n+1}\longrightarrow\OO_{C_{n+1}}\longrightarrow\OO_{C_n}\longrightarrow 0
\end{equation*}
（即 $C_n$ 在 $C_{n+1}$ 中的\trClosedSubscheme{}正合列 \eqref{eq:closed-subscheme-exact-sequence}），并利用典范\trIsomorphism{} $\mathrm{H}^0(\mathbb{P}^1,\OO_{\mathbb{P}^1}(n))\stackrel{\sim}{\hookrightarrow}\Sym^n\bigl(\mathrm{H}^0(\mathbb{P}^1,\OO_{\mathbb{P}^1}(1))\bigr)$。棘手之处在于：我们要的这个同构是作为\emph{环}的同构。
\textbf{28.6.F.（练习）}\quad 证明这一点。（这可能让你回想起，我们曾在\mbox{练习~19.11.E} 中求出一个\trInvertibleSheaf{}的\trTotalSpace{}上的函数环。这并非巧合。）提示：把 $\mathscr{I}^n/\mathscr{I}^{n+1}$ 解读为 $\Sym^n(\mathscr{I}/\mathscr{I}^2)$（\mbox{推论~22.3.9}），进而把它解读为 $\OO_{\mathbb{P}^1}(n)$，并证明：取整体\trSection{}之后，\eqref{eq:castelnuovo-graded-piece-seq} 仍然正合。到这里，你将已经证明：$\mathrm{H}^0(C_n,\OO_{C_n})$ 的某种\trAssociatedGraded{}环同构于 $k[[x,y]]/(x,y)^n$。请证明这就足以建立所欲求的环同构。
于是我们得到了\trSmooth{}性，第 2 步完成。
\emph{第 3 步.}\quad 现在我们必须认出 $\pi$ 就是 $Y$ 在 $q$ 处的\trBlowup。由于 $C$ 是 $X$ 上的一个\trEffectiveCartierDivisor，由胀开的\trUniversalProperty{}可知\footnote{原书此处作「by the universal property of blowing up means that」，「by $\ldots$」与「means that」两句式杂糅，兹按文义译出。——译者注}，存在唯一的\trMorphism{} $\alpha$ 使下图交换：
\[
\begin{tikzcd}[row sep=2.8em, column sep=4.5em]
X \ar[rr, "\alpha"] \ar[dr, "\pi"'] && \mathrm{Bl}_q\,Y \ar[dl, "\beta"] \\
& Y. &
\end{tikzcd}
\]
由针对\trProjective\trMorphism{}的\trCancellationTheorem{}（\mbox{练习~17.3.C}），$\alpha$ 是\trProjective{}的。
\textbf{28.6.G.（练习）}\quad 证明：$\alpha$ 具有\trFinite\trFiber。
由于具有有限\trFiber{}的\trProjective\trMorphism{}是\trFinite{}的（\mbox{定理~18.1.6}），$\alpha$ 是一个\trFinite\trMorphism。但这样一来，$\alpha$ 是一个到\trNormal\trScheme{}的\trBirational\trFinite\trMorphism，从而由\mbox{练习~10.7.P} 是一个\trIsomorphism。\hfill$\square$
到目前为止，我们只把\mbox{形式函数定理~28.3.2} 用在了 $\mathrm{H}^0$ 上。下一道练习将给你一个机会，看看它在其他\trCohomology{}群上的实际运作。
\textbf{28.6.H.（练习）}

(a) 设 $\beta\colon X=\mathrm{Bl}_q\,Y\to Y$ 是 $k$ 上的一张\trSmooth\trSurface{}在一个\trSmooth{}的 $k$-\trValuedPoint{}处的\trBlowup。证明：$R^1\beta_*\OO_X=R^2\beta_*\OO_X=0$。

(b) 证明：自然映射 $\beta^*\colon\mathrm{H}^1(Y,\OO_Y)\to\mathrm{H}^1(X,\OO_X)$ 与 $\beta^*\colon\mathrm{H}^2(Y,\OO_Y)\to\mathrm{H}^2(X,\OO_X)$ 都是\trIsomorphism。

\textbf{28.6.2. 初等变换。}

回忆\mbox{例~17.2.6} 中给出的\trHirzebruchSurface{} $\mathbb{F}_n$（域 $k$ 上的）的定义。若 $n>0$，则 $\mathbb{F}_n$ 上有唯一一条\trSelfIntersection{}最负的\trCurve{} $E$（\mbox{练习~20.2.R}）。它的\trSelfIntersection{}为 $E\cdot E:=\deg_E\OO(E)|_E=-n$。它是投射 $\pi\colon\mathbb{F}_n\to\mathbb{P}^1$ 的一个\trSection{}。（若 $n=0$，此时 $\mathbb{F}_n\cong\mathbb{P}^1\times\mathbb{P}^1$，则 $E$ 通常被定义为 $\pi\colon\mathbb{P}^1\times\mathbb{P}^1\to\mathbb{P}^1$ 的任意一个「常值」\trSection{}——它并不唯一。）
设 $q$ 是 $\mathbb{F}_n$ 的一个 $k$-\trValuedPoint。令 $F$ 为 $\pi$ 的含有 $q$ 的那条\trFiber。
\textbf{28.6.I.（练习）}\quad 证明：$F\cdot F:=\deg_F\OO(F)|_F=0$。
设 $\beta\colon\mathrm{Bl}_q\,\mathbb{F}_n\to\mathbb{F}_n$ 为在 $q$ 处的\trBlowup，并令 $Z$ 为其\trExceptionalDivisor（因为「$E$」这个名字已被占用）。令 $F'$ 为 $F$ 的\trProperTransform，$E'$ 为 $E$ 的\trProperTransform。
\textbf{28.6.J.（练习）}\quad 证明：$F'\cdot F':=\deg_{F'}\OO(F')|_{F'}=-1$。
于是，由\trCastelnuovo{}判别准则，我们可以把 $F'$ 收缩掉，从而得到一张新的\trSurface{} $Y$。
\textbf{28.6.K.（练习）}\quad 设 $n>0$。若 $q\notin E$，证明：$Y\cong\mathbb{F}_{n-1}$。若 $q\in E$，证明：$Y\cong\mathbb{F}_{n+1}$。可能的提示：倘若你已经知道 $Y$ 是\trHirzebruchSurface{} $\mathbb{F}_m$，那么当 $m>0$ 时，你可以凭借那条唯一具有负\trSelfIntersection{}的\trCurve{}的\trSelfIntersection{}把它认出来。
\textbf{28.6.L.（练习）}\quad 设 $n=0$。证明：$Y\cong\mathbb{F}_1$。提示：\mbox{练习~22.4.K}。
这里的讨论可以推广到更一般的\trCurve{}上的 $\mathbb{P}^1$-丛。下面不过是这个故事的第一步。
\textbf{28.6.M.（练习）}\quad 设 $C$ 是 $k$ 上一条\trSmooth\trCurve，$\pi\colon X\to C$ 是 $C$ 上的一个 $\mathbb{P}^1$-丛。若 $q$ 是 $X$ 的一个 $k$-\trValuedPoint，那么我们可以胀开 $q$，再收缩经过 $q$ 的那条\trFiber{}的\trProperTransform。证明：其结果仍是 $C$ 上的一个 $\mathbb{P}^1$-丛 $X'$。（这称为该\trRuling{}曲面的一个\textbf{初等变换}（elementary transformation）。）
\textbf{28.6.N.（练习）}\quad 设 $X=\mathit{Proj}_C(\mathscr{L}\oplus\mathscr{M})$，而 $q$ 落在相应于正合列
\[
0\longrightarrow\mathscr{L}\longrightarrow\mathscr{L}\oplus\mathscr{M}\longrightarrow\mathscr{M}\longrightarrow 0.
\]
的\trSection{}之上。证明：$X'\cong\mathit{Proj}_C(\mathscr{L}(\pi(q))\oplus\mathscr{M})$。\footnote{原书此处作「$\mathit{Proj}_C(\mathscr{L}(\pi(q))\oplus\mathscr{M}))$」，末尾多出一个右括号，径删。——译者注}
倘若只告诉你 $X\cong\mathit{Proj}_C\,\mathscr{V}$（其中 $\mathscr{V}$ 是一个秩为 $2$ 的\trLocallyFreeSheaf），而 $q$ 落在相应于 \eqref{eq:section-filtration-sequence} 的\trSection{}之上，你能猜测（或者证明）会发生什么吗？

\textbf{28.6.3. 曲面的极小模型。}

代数闭域上的一张\trSurface，若不含 $(-1)$-\trCurve，就称为\textbf{极小的}（minimal）。若\trSurface{} $X$ 不是极小的，我们可以选取一条 $(-1)$-\trCurve{}并把它收缩。若所得\trSurface{}仍不是极小的，我们可以再收缩一条 $(-1)$-\trCurve，如此等等。由\trPicard{}数 $\rho$ 的有限性（这一点我们诚然未曾证明；见 \S18.4.12），再加上收缩会使 $\rho$ 减小 $1$（据\mbox{练习~22.4.O}），这个过程必定终止。由此我们构造出 $X$ 的一个\textbf{极小模型}（minimal model）——一张\trBirational{}于 $X$ 的极小\trSurface。倘若极小模型是唯一的，这便提供了一种按双有理等价给\trSurface{}分类的办法：若 $X$ 与 $X'$ 是\trSmooth\trProjective{}\trSurface，想知道它们是否\trBirational，我们只须找出它们各自的极小模型，再看它们是否\trIsomorphic。
遗憾的是，极小模型并不唯一。$\mathbb{P}^1\times\mathbb{P}^1$ 与 $\mathbb{P}^2$ 的例子表明，同一个双有理等价类中可以有多于一个极小模型。（为什么这两张曲面都是极小的？）此外，$\mathbb{P}^2$ 在两点处的\trBlowup{}与 $\mathbb{P}^1\times\mathbb{P}^1$ 在一点处的\trBlowup{}之间的\trIsomorphism{}（\mbox{练习~22.4.K}）表明，一张\trSmooth\trSurface{}可以拥有多于一个极小模型。
尽管如此，这种不唯一性已被理解得相当透彻：\trZariski{}曾证明，具有多于一个极小模型的\trSurface{}恰是那些\trBirational{}于 $C\times\mathbb{P}^1$（$C$ 为某条\trCurve）的曲面。这是曲面的\trEnriques--\trKodaira{}分类的一部分。欲知更多，见 [Be]、[Ba] 或 [BPV, \S 6]。
\textbf{28.6.4.}\quad 这里还有两条值得提及但在此不拟证明的事实。其一，奇异的\trSurface{}可以先作\trNormalization、再反复胀开点，从而\trDesingularization；这一显式的\trDesingularization{}与 \S28.4.4 中对\trCurve{}的那个版本相类似，并且被\trHironaka{}的\trResolutionOfSingularities{}定理（\mbox{注记~22.4.6}）所推广。其二，\trSmooth\trSurface{}之间的\trBirational\trRationalMap{} $X\leftrightarrow Y$ 可以分解为有限多次胀开接以有限多次收缩（参看 \S22.4.12，论及弱分解）。作为推论，据\mbox{练习~28.6.H}，\trBirational{}的\trProjective{}\trSurface{}具有相同的 $\mathrm{H}^i(\OO)$（$i=0$、$1$、$2$）。（利用 \S21.5.10 中陈述的事实，你还可以证明：在 $\mathbb{C}$ 上，这意味着除 $h^{1,1}$ 之外的整个\trHodge{}菱形都是\trSurface{}的双有理不变量。）
